\chapter{Plasticity}

We first focus on plasticity in one dimension, 
as opposed to fully three-dimensional plasticity, 
to emphasize concepts while being
as mathematically simple as possible.  
After the concepts are clear, the extension of the mathematics from
one- to three-dimensions can be made naturally.  

\section{Perfect Plasticity}
The extension of elasticity to plasticity is made by way of 
{\em perfect} plasticity.
Conceptually, perfect plasticity in the constitutive stress-strain 
relationship is obtained by ``capping'' the
elastic response function at some yield stress 
(to be defined later as $\sigma_y$).  Beyond the yield stress, the 
stress response function is perfectly horizontal.  
Increases in strain produce no increase in stress
beyond the yield stress.  


\subsection{Deformation Measures}
Let the body $\Omega$ have strain measure 
$\hstrain : \Omega \rightarrow \real$ 
in one dimension be defined as a ratio of change in length, that is 
the difference between the new (deformed, current configuration) 
length $\ell$ of body $\Omega$ and the original (undeformed, reference 
configuration) length $L$, divided by the original length, 
\be
 \hstrain \defe \frac{\ell - L}{L}.
 \ee
 The total strain $\hstrain$ in the body be defined as an 
 additive decomposition of elastic strain $\hstraine$ and 
plastic strain $\hstrainp$, 
\be
 \label{eq:smallStrainAddition}
 \hstrain \defe \hstraine + \hstrainp.
\ee


\subsection{Constitutive Law}
Then, the traditional stress-strain relationship, $\sigma = E \hstraine$, 
can be rewritten as
\be
 \label{eq:stressStrain}
 \sigma = E \left( \hstrain - \hstrainp \right).
\ee
A change in the configuration of the plastic element 
(see  Fig.~\ref{fig:plastic_slider})
can only occur if the notion of {\bf slip rate} is
allowed.

%\begin{figure}[htb]
% \psfrag{a}{$\sigma$}
% \psfrag{b}{$\sigma_y$}
% %\fbox{\begin{minipage}{0.95\columnwidth}\centering
% \begin{center}
% \includegraphics[width=2.8in]{fig/plasticityFrictionElement.eps}
% \end{center}
% \caption{Phenomenological model of a friction element.  A unit cross-sectional
% area is assumed.  Motion in the friction 
% element occurs only if the friction force, $\sigma_y$, 
% is exceeded.  The motion, in rate form, is
% referred to as the slip rate. Note that we employ a slight abuse of 
% notation by using $\sigma$ to describe a force $F$, since the cross-sectional area $A$ is
% unity, $F = \sigma A \implies \sigma \cdot 1 = \sigma$.}
% \label{fig:plasticFrictionElement_old} % label must come after caption 
% %\end{minipage}
% %} % \fbox
%\end{figure}


\begin{figure}[htb]
  \begin{center}
  %\begin{overpic}[width=0.5\textwidth, grid, tics=10]{fig/plastic_slider}
  %\begin{overpic}[width=0.5\textwidth,natwidth=144,natheight=51]{fig/plastic_slider.pdf}
  \begin{overpic}[width=0.5\textwidth]{fig/plastic_slider.pdf}
    \put( 4, 13){$\sigma$}
    \put(95, 13){$\sigma$}
    \put(47, 23){$\sigma_y$}
  \end{overpic}
  \end{center}
 \caption{Phenomenological model of a friction element.  A unit cross-sectional
 area is assumed.  Motion in the friction 
 element occurs only if the friction force, $\sigma_y$, 
 is exceeded.  The motion, in rate form, is
 referred to as the slip rate. Note that we employ a slight abuse of 
 notation by using $\sigma$ to describe a force $F$, since the 
 cross-sectional area $A$ is
 unity, $F = \sigma A \implies \sigma \cdot 1 = \sigma$.}
 \label{fig:plastic_slider} % label must come after caption 
\end{figure}

%\begin{figure}[htb]
%\begin{center}
%  \psfragfig[width=6cm]{fig/plasticityFrictionElement}{% the .eps file
%    \psfrag{a}{$\sigma$}
%    \psfrag{b}{$\sigma_y$}
%  } % psfragfig
%\end{center}
% \caption{Phenomenological model of a friction element.  A unit cross-sectional
% area is assumed.  Motion in the friction 
% element occurs only if the friction force, $\sigma_y$, 
% is exceeded.  The motion, in rate form, is
% referred to as the slip rate. Note that we employ a slight abuse of 
% notation by using $\sigma$ to describe a force $F$, since the cross-sectional area $A$ is
% unity, $F = \sigma A \implies \sigma \cdot 1 = \sigma$.}
% \label{fig:plasticFrictionElement} % label must come after caption 
%\end{figure}
%


That is, {\em assume} that $\hstrainp(t)$ is a 
function of time on interval $[0, T] \subset \real$.  
Then from~(\ref{eq:smallStrainAddition}) and~(\ref{eq:stressStrain}), 
time evolution expressions $\hstrain(t)$ and $\sigma(t)$ exist, respectively.  
Finally, the presence of evolution equations allows the constitutive law 
in~(\ref{eq:stressStrain}) to be
written as
\be
 \sigmadot = E \left( \hstraindot - \hstrainpdot \right).
 \label{eq:stressStrainRate}
\ee
This form of the constitutive equation will 
become useful for quantifying plastic 
flow, as seen below in~(\ref{eq:PhiChainRule}).

%\begin{figure}[htb]

%%\begin{figure*}[h!]
% \psfrag{a}{(b) slip rate $\gammadot \geq 0$}
% \psfrag{b}{(a) yield function $\Phi(\sigma) \leq 0$}
% \psfrag{c}{(c) consistency}
% \psfrag{g}{$\gammadot$}  
% \psfrag{p}{$\Phi$} 
% %\fbox{\begin{minipage}{0.98\textwidth}\centering
% \begin{center}
% \includegraphics[width=5.5in]{fig/plasticunion.eps}
% \end{center}
% \caption{Unilateral conditions (a) $\Phi(\sigma) \leq 0$ and 
% (b) slip rate $\gammadot \geq 0$.  
%  The complementarity condition, $\gammadot \; \Phi(\sigma) = 0$, 
%  is shown in~(c).}
% \label{fig:plasticunion} % label must come after caption 
% %\end{minipage}
% %} % \fbox
%\end{figure}
%%\end{figure*}
%
%
%\begin{figure}[htb]
%\begin{center}
%  \psfragfig[width=\textwidth]{fig/plasticunion}{% the .eps file
%    \psfrag{a}{(b) slip rate $\gammadot \geq 0$}
%    \psfrag{b}{(a) yield function $\Phi(\sigma) \leq 0$}
%    \psfrag{c}{(c) consistency}
%    \psfrag{g}{$\gammadot$}  
%    \psfrag{p}{$\Phi$} 
%  } % psfragfig
%\end{center}
% \caption{Unilateral conditions (a) $\Phi(\sigma) \leq 0$ and 
% (b) slip rate $\gammadot \geq 0$.  
%  The complementarity condition, $\gammadot \; \Phi(\sigma) = 0$, 
%  is shown in~(c).}
% \label{fig:plasticunion_new} % label must come after caption 
%\end{figure}
%

\subsection{Yield Function}
The stress $\sigma$ in the frictional device can never greater in 
absolute value than the yield stress (also known
as flow stress) $\sigma_y \in \realplus$.
The {\bf yield function} $\Phi$ is then defined as 
\be
 \Phi(\sigma) \defe \abs{\sigma} - \sigma_y \leq 0.
 \label{eq:yieldFunctionPerfectPlastic}
\ee
For stress levels within the elastic domain, 
only elastic straining can occur, thus
\be
 \mbox{If} \;\;\; \Phi(\sigma) < 0 \Rightarrow \hstrainpdot = 0.
\ee
For stress levels on the boundary of the elastic domain 
(at the yield stress), either plastic strain via plastic loading or
non-plastic strain via elastic unloading occurs.
\begin{align}
 \mbox{If} \;\; \Phi(\sigma) & = 0  \nonumber \\
  & \Rightarrow  
   \left\{
     \begin{tabular}{ll}
     $\hstrainpdot \neq 0$ & for plastic loading, \\
     $\hstrainpdot = 0$ & for elastic unloading.
     \end{tabular}
   \right.
\end{align}

\subsection{Plastic Flow Rule, Loading/Unloading Conditions}
We next define how, on the yield surface, the plastic strain evolves 
through time, which is called the 
{\bf plastic flow rule}.  A reasonable first approximation for the 
plastic strain evolution is to posit
\be
 \hstrainpdot \defe \gammadot \; \sign (\sigma),
 \label{eq:flowrule}
\ee
where $\gammadot \geq 0$ is the {\bf plastic multiplier} or 
absolute value of the plastic slip rate, and 
where $\sign : \real \rightarrow \{-1,+1\}$ is the {\em signum} 
function defined as
\be
 \sign(a) = 
 \left\{ 
   \begin{tabular}{rl}
    $+1$ & if $a \geq 0$, \\
    $-1$ & if $a < 0$.
   \end{tabular}
 \right.
\ee
Both $\sigma$ and $\gammadot$ are restricted by {\bf unilateral constraints}, 
\be
 \begin{tabular}{r}
  $\Phi(\sigma) \leq 0$ \\
   $\gammadot \geq 0$ 
 \end{tabular}
\ee
Now if $\Phi(\sigma) < 0$, then $\gammadot = 0$.  
Similarly, if $\gammadot > 0$, then $\Phi(\sigma) = 0$.
These two observations require
\be
\gammadot \; \Phi(\sigma) = 0,
\ee
which is referred to as the {\bf complementarity condition}, 
and when taken with the unilateral constraints, describes the
so-called {\bf Kuhn-Tucker conditions}.
The unilateral constraints and the complementarity condition are illustrated 
in Fig.~\ref{fig:plastic_union}.

\begin{figure}[htb]
  \begin{center}
  %\begin{overpic}[width=\textwidth, grid, tics=10]{fig/plastic_union}
  %\begin{overpic}[width=0.95\textwidth,natwidth=414,natheight=157]{fig/plastic_union.pdf}
  \begin{overpic}[width=0.95\textwidth]{fig/plastic_union.pdf}
    \put(28, 13){$\gammadot$}
    \put(63, 13){$\gammadot$}
    \put(99, 13){$\gammadot$}
    \put(12.5, 29){$\Phi$}
    \put(47, 29){$\Phi$}
    \put(82, 29){$\Phi$}
    \put( 0, 35){(a) yield function $\Phi(\sigma) \leq 0$}
    \put(39, 35){(b) slip rate $\gammadot \geq 0$}
    \put(74, 35){(c) consistency}
  \end{overpic}
  \end{center}
 \caption{Unilateral conditions (a) $\Phi(\sigma) \leq 0$ and 
 (b) slip rate $\gammadot \geq 0$.  
  The complementarity condition, $\gammadot \; \Phi(\sigma) = 0$, 
  is shown in~(c).}
 \label{fig:plastic_union} % label must come after caption 
\end{figure}


Together, the unilateral constraints and complementarity condition 
create the normal plasticity law, shown in 
Fig.~\ref{fig:plastic_law}.

\begin{figure}[htb]
  \begin{center}
  %\begin{overpic}[width=0.3\textwidth, grid, tics=10]{fig/plastic_law}
  %\begin{overpic}[width=0.3\textwidth,natwidth=126,natheight=130]{fig/plastic_law.pdf}
  \begin{overpic}[width=0.3\textwidth]{fig/plastic_law.pdf}
    \put(90, 41){$\gammadot$}
    \put(41, 90){$\Phi$}
  \end{overpic}
  \end{center}
  \caption{The normal plastic law described as the intersection of the 
  unilateral constraints and the complementarity condition.}
 \label{fig:plastic_law} % label must come after caption 
\end{figure}

%\begin{figure}[htb]
% \psfrag{g}{$\gammadot$}  
% \psfrag{p}{$\Phi$} 
% %\fbox{\begin{minipage}{0.95\columnwidth}\centering
%  \includegraphics[width=1.8in]{fig/plasticlaw.eps}
%  \caption{The normal plastic law described as the intersection of the 
%  unilateral constraints and the complementarity condition.}
%  \label{fig:plasticlaw} % label must come after caption 
%  %\end{minipage}
%  %} % fbox
%\end{figure}

Next, consider any time interval $[t, t + \Delta t]$ 
during which {\em plastic flow occurs}, the yield function 
remains constant, $\Phi(\sigma) = 0$.  Similarly, 
if during this time interval, 
we have elastic response only, and $\Phidot(\sigma) < 0$, 
then necessarily $\gammadot = 0$.  These relationships create 
the {\bf persistency (or consistency) condition}
\be
 \left.
  \begin{tabular}{l}
   $\mbox{If} \;\;\; \gammadot > 0 \Rightarrow \Phidot = 0$, \\
   $\mbox{If} \;\;\; \Phidot < 0 \Rightarrow \gammadot = 0$,
  \end{tabular}
  \right\} 
  \;\;\; \Rightarrow \;\;\;
  \gammadot \; \Phidot = 0.
 \label{eq:persistencyCondition}
\ee



\subsection{Determination of Plastic Multiplier}
As of yet, the value of the plastic multiplier $\gammadot$ 
has remained undetermined.  For non-trivial $\gammadot$, we
know $\Phidot = 0$ by the persistency condition.  
An explicit form for $\Phidot$ may be found using 
the chain rule with~(\ref{eq:stressStrainRate}) 
and~(\ref{eq:flowrule}), 
\begin{align}
 \frac{\partial \Phi}{\partial t} &= \dst\frac{\partial \Phi}{\partial \sigma} 
 \dst\frac{\partial \sigma}{\partial t} 
 \nonumber \\
 &=  \dst\frac{\partial }{\partial \sigma} \left(\abs{\sigma} - 
 \sigma_y \right) E \left( \hstraindot - \hstrainpdot \right) \nonumber \\
 &= \sign(\sigma) E \left( \hstraindot - \gammadot \; \sign(\sigma) \right), 
 \label{eq:PhiChainRule}
\end{align}
or, compactly
\be
 \Phidot = \sign(\sigma) E \hstraindot - E \gammadot \; \setequal \;  0,
\ee
gives
\be
 \gammadot = \hstraindot \; \sign(\sigma)
 \label{eq:gammaDotPerfectPlasticity}
\ee
which, when compared with~(\ref{eq:flowrule}), shows that for 
perfect plasticity during plastic flow, 
\be
 \hstrain = \hstrainp.
 \label{eq:strainPerfectPlasticity}
\ee
The results from~(\ref{eq:gammaDotPerfectPlasticity}) 
and~(\ref{eq:strainPerfectPlasticity}) 
are remarkable for the following reasons:  During plastic flow, 
\begin{itemize}
 \item None of the strain $\hstrain$ is {\em elastic} strain $\hstraine$; 
 100\% of the stain $\hstrain$ is {\em plastic} strain $\hstrainp$.
 This feature will not be true, in general, for more sophisticated yield 
 functions (for example, that have plastic hardening).
 \item The plastic flow parameter $\gammadot$ is equal to the strain 
 $\hstrain$ times the sign of the stress. 
 For more sophisticated yield functions, the plastic flow 
 parameter will not by identically 
 equal to the signed strain, but will be related thereto via 
 factors of the elastic and hardening parameters, as shown subsequent
 developments and in particular in~(\ref{eq:gammaDotHardening}).
\end{itemize}

\section{Hardening Plasticity}

Hardening plasticity is considered an extension of perfect plasticity.  With 
perfect plasticity, the slope of stress-strain relationship at yielding is
identically zero.  With hardening plasticity, the slope becomes 
positive.\footnote{For completeness, we note that softening plasticity occurs
when the slope becomes negative.}   

\subsection{Remarks on Loss of Stability}
Perfect plasticity allows for the notion of unbounded 
strains when stress is at the yield stress. 
The slope of the stress-strain curve for strains 
beyond the elastic strain region is identically zero.
This zero slope causes numerical problems for the 
consistent tangent, which is inverted and multiplied by
the external load to obtain displacements.  
Three strategies to circumvent this zero tangent are 
mentioned here:  arc-length, displacement control, and dynamic formulation.

The most notable of the continuation methods is called the  
{\bf arc-length method}, 
which is a non-trivial area of study within computational mechanics.  
Arc-length methods will be used eventually, but their
explication is delayed for now, in favor development of hardening plasticity.  

Despite being more slightly more complex than
perfect plasticity, hardening plasticity does not 
likewise suffer from loss of ellipticity and therefore is actually easier to solve since
algorithmic enhancements, such as arc-length, are 
(usually\footnote{If the hardening is ``soft'' enough that the hardening 
approximates perfect plasticity, hardening plasticity can 
result in lack of conditioning.  Usually, however, 
hardening is sufficiently stiff and
therefore loss of ellipticity is not a concern.}) unnecessary.

Next, algorithmic formulations that utilize {\bf displacement control} 
instead of force control can avoid loss of ellipticity if the 
displacement boundary conditions reduce the number of equations 
in the tangent to provide rank sufficiency.  
However, displacement control is less-favored than 
force control since we often know imposed forces and 
seek to discover the displacement response.  The setting of 
known inhomogeneous displacements and unknown forces is
less frequently encountered in practice.  
On the topic of load control versus displacement 
control, de~Borst~{\em et al.,}\footnote{De Borst, R., 
Crisfield, M. A., Remmers, J. J., \& Verhoosel, C. V. (2012). 
Nonlinear finite element analysis of solids and structures. John Wiley \& Sons, at page 50--51.} offer similar insights and additional context. 

Finally, elastoplastic {\bf dynamics}, and so-called dynamic 
relaxation as a subset thereof,
are additional strategies to circumvent the above-mentioned 
problem as these approaches condition the
consistent tangent with mass matrices.  The benefit of 
mass matrices comes at the cost of increased expense to compute 
and store the same, and to incorporate
a time stepping algorithm, such as the Newmark methods.

\subsection{Yield Function}
The departure of hardening plasticity from perfect plasticity 
begins with the yield function.
The {\bf yield function} from perfect 
plasticity in~(\ref{eq:yieldFunctionPerfectPlastic}) 
is now modified to include a {\bf linear hardening modulus} $H \in \real$ and
the {\bf internal hardening variable} $\alpha : [0, T] \rightarrow \real$,
\be
 \Phi(\sigma, \hstrainp) \defe \abs{\sigma} - \left( \sigma_y + H \alpha \right) \leq 0.
 \label{eq:yieldFunctionHardening}
\ee
for strain-hardening plasticity.
The variable $\alpha : [0, T] \rightarrow \real$ is called the 
{\bf internal hardening variable} and is a nonnegative function of the amount of 
plastic flow (slip).  
The simplest evolution for $\alpha$ can be written as
\be
 \alphadot = \abs{\hstrainpdot}.
\ee
The {\bf plastic flow rule}, 
\be 
 \hstrainpdot = \gammadot \; \sign(\sigma),
 \label{eq:plasticFlowRuleHarding}
\ee
remains unchanged.  The {\bf Kuhn-Tucker conditions} are rewritten to include 
dependence on $\alpha$, 
\be
  \gammadot \geq 0 \;\;\;\;\; 
  \Phi(\sigma, \alpha) \leq 0 \;\;\;\;\;
  \gammadot \; \Phi(\sigma, \alpha) = 0, 
\ee
and the {\bf persistency condition} is written as
\be
 \gammadot \; \Phidot(\sigma, \alpha) = 0.
\ee

\subsection{Determination of Plastic Multiplier}
Like perfect plasticity, strain-hardening plasticity uses the time derivative
of the yield function to generate an expression for the plastic multiplier
$\gammadot$, which ultimately leads to a relationship between stress rates and
strain rates, and the elasto-plastic tangent.

With the yield function in~(\ref{eq:yieldFunctionHardening}),
\begin{align}
 \frac{\partial \Phi}{\partial t} 
  &=  \dst\frac{\partial \Phi}{\partial \sigma} 
  \dst\frac{\partial \sigma}{\partial t} 
  +\dst\frac{\partial \Phi}{\partial \alpha} 
  \dst\frac{\partial \alpha}{\partial t}, \nonumber \nonumber \\
  &= \sign(\sigma) \sigmadot - H \alphadot, \nonumber \\
  &= \sign(\sigma) E \left( \hstraindot - \hstrainpdot \right) 
  - H \abs{\hstrainpdot}, \nonumber \\
  &= \sign(\sigma) E \hstraindot - E \gammadot - H \gammadot \; \setequal \;  0, 
\end{align}
gives
\be
 \boxed{\gammadot = \hstraindot \; \sign(\sigma) \dst\frac{E}{E + H} }.
 \label{eq:gammaDotHardening}
\ee
The ratio $\frac{E}{E + H}$ in~({\ref{eq:gammaDotHardening}) 
may be conceived as the apportionment of
total strain (rate) to plastic slip:
\begin{itemize}
 \item If $H = 0$, then all strain is plastic slip, 
 and the perfect plasticity results obtained in previous sections 
 are obtained as a special case.
 \item If $H=E$, then half of all strain goes into plastic slip.
 \item If $H \gg E$, then very little $\hstraindot$ goes into plastic slip.
\end{itemize}

\subsection{Elasto-Plastic Tangent}
The elasto-plastic tangent is found by 
observing $\sigmadot = \frac{d\sigma}{d\hstrain} \hstraindot$.  
Starting with the elastic stress-strain relationship
$\sigma = E \left(\hstrain - \hstrainp\right)$,
in the rate form and 
substituting for $\hstraindot$ from~(\ref{eq:gammaDotHardening})
and for $\hstrainpdot$ from the plastic flow 
rule~(\ref{eq:plasticFlowRuleHarding}) 
gives
\begin{align}
 \sigmadot 
 & = E \left( \hstraindot - \hstrainpdot \right), \nonumber \\
 & = \gammadot \; \sign(\sigma) (E + H) - \gammadot \; \sign(\sigma) E, \nonumber \\ 
 & = \gammadot \; \sign(\sigma) H.
\end{align}
Substituting for $\gammadot$ from~(\ref{eq:gammaDotHardening}) gives
\be
 \sigmadot = \dst\frac{E H}{E + H} \hstraindot, 
\ee
where $\Eep$, the elasto-plastic tangent, is given by
\be
 \boxed{\Eep = \dst\frac{E H}{E + H}}.
\ee
		      
The concept of the elasto-plastic tangent is readily understood with 
the phenomenological model
shown in Fig.~\ref{fig:plastic_slider_harden}.


%\begin{figure}[htb]
% \psfrag{a}{$\sigma$}
% \psfrag{b}{$\sigma_y$}
% \psfrag{c}{$H$}
% \psfrag{d}{$E$}
% \psfrag{e}{$\hstrainp$}
% \psfrag{f}{$\hstraine$}
% \psfrag{g}{$\hstrain$}
% %\fbox{\begin{minipage}{0.98\columnwidth}\centering
% \begin{center}
% \includegraphics[width=3.0in]{fig/plasticitySpringModel.eps}
% \end{center}
% \caption{Phenomenological model of strain-hardening 
% plasticity with applied stress $\sigma$, yield stress $\sigma_y$, 
% hardening modulus $H$ and elastic modulus $E$.  The 
% total strain, $\hstrain$ is the sum of plastic strain $\hstrainp$ 
% and elastic strain $\hstraine$.}
% \label{fig:plasticSpringModel} % label must come after caption 
% %\end{minipage}
% %} % fbox
%\end{figure}

\begin{figure}[htb]
  \begin{center}
  %\begin{overpic}[width=0.6\textwidth, grid, tics=10]{fig/plastic_slider_harden}
  %\begin{overpic}[width=0.5\textwidth,natwidth=216,natheight=195]{fig/plastic_slider_harden.pdf}
  \begin{overpic}[width=0.5\textwidth]{fig/plastic_slider_harden.pdf}
    \put(66, 72){$E$}
    \put(32, 54){$H$}
    \put( 1, 61){$\sigma$}
    \put(97, 61){$\sigma$}
    \put(32, 84){$\sigma_y$}
    \put(32, 28){$\hstrainp$}
    \put(66, 28){$\hstraine$}
    \put(49, 10){$\hstrain$}
  \end{overpic}
  \end{center}
 \caption{Phenomenological model of strain-hardening 
 plasticity with applied stress $\sigma$, yield stress $\sigma_y$, 
 hardening modulus $H$ and elastic modulus $E$.  The 
 total strain, $\hstrain$ is the sum of plastic strain $\hstrainp$ 
 and elastic strain $\hstraine$.}
 \label{fig:plastic_slider_harden} % label must come after caption 
\end{figure}


The stress $\sigma$ is applied to the model.  
Initially, the applied stress $\sigma$ creates only
elastic strain $\hstraine$ since $\sigma \leq \sigma_y$.  
When $\sigma > \sigma_y$, strain composed
of {\em both} elastic strain $\hstraine$ and plastic 
strain $\hstrainp$ is created.  The linear moduli $E$ and $H$ 
act as springs in series.  The relationship for two springs 
with spring constants, $k_1$ and $k_2$, in series immediately shows the
relationship between the equivalent spring constant $\keq$ 
and the elasto-plastic tangent $\Eep$
\begin{align}
 \frac{1}{\keq} = \frac{1}{k_1} + \frac{1}{k_2} 
 \implies & \keq = \frac{k_1 k_2}{k_1 + k_2} \\
 & \; \Updownarrow \hspace{1.1cm} \Updownarrow \nonumber \\
 & \Eep = \frac{E H}{E + H}
\end{align}
The hardening modulus $H$ is shown in Fig.~\ref{fig:stress_strain_harden}(b)
such that, during plastic flow, the incremental stress $\partial \sigma$ relative
to the incremental plastic strain $\partial\hstrainp$ is given by
\be
 \frac{\partial \sigma}{\partial \hstrainp} = H.
\ee

%%\begin{figure*}
%\begin{figure}[htb]
% \psfrag{1}{(a)}
% \psfrag{2}{(b)}
%  \psfrag{a}{elastic}
%  \psfrag{b}{elasto-plastic}
%  \psfrag{c}{plastic}
%  \psfrag{E}{$E$}
%  \psfrag{e}{$\hstrain$}
%  \psfrag{F}{$\Eep$}
%  \psfrag{f}{$\hstrainp$}
%  \psfrag{G}{$H$}
%  \psfrag{s}{$\sigma$}
%  \psfrag{y}{$\sigma_y$}
%  %\fbox{\begin{minipage}{0.98\textwidth}\centering
%  \includegraphics[width=\textwidth]{fig/plasticityHardening.eps}
%  \caption{(a) Strain-hardening stress-strain relationship with elastic 
%  and elasto-plastic strain regions delineated, elastic modulus $E$, 
%  elasto-plastic modulus $\Eep$, and yield stress $\sigma_y$, and 
%  (b) stress-(plastic) strain relationship with hardening modulus $H$.}
%  \label{fig:plasticityHardening} % label must come after caption 
%  %\end{minipage}
%  %} % fbox
%%\end{figure*}
%\end{figure}


\begin{figure}[htb]
  \begin{center}
  %\begin{overpic}[width=\textwidth, grid, tics=10]{fig/stress_strain_harden}
  %\begin{overpic}[width=\textwidth,natwidth=504,natheight=288]{fig/stress_strain_harden.pdf}
  \begin{overpic}[width=\textwidth]{fig/stress_strain_harden.pdf}
  \put( 7.25, 10){elastic}
  \put(20, 10){elasto-plastic}
  \put(72, 10){plastic}
  \put(11, 17){$E$}
  \put(45, 14){$\hstrain$}
  \put(32, 36){$\Eep$}
  \put(94, 14){$\hstrainp$}
  \put(73, 36){$H$}
  \put( 5, 52){$\sigma$}
  \put(55, 52){$\sigma$}
  \put( 4, 30){$\sigma_y$}
  \put(54, 30){$\sigma_y$}
  \put( 0, 55){(a)}
  \put(52, 55){(b)}
  \end{overpic}
  \end{center}
  \caption{(a) Strain-hardening stress-strain relationship with elastic 
  and elasto-plastic strain regions delineated, elastic modulus $E$, 
  elasto-plastic modulus $\Eep$, and yield stress $\sigma_y$, and 
  (b) stress-(plastic) strain relationship with hardening modulus $H$.}
 \label{fig:stress_strain_harden} % label must come after caption
\end{figure}


\section{Elasto-Plastic 2D Truss Element}

Consider the two-dimensional truss member shown in 
Fig.~\ref{fig:nonlinear_truss}.
The coordinates in the reference configuration are noted 
with $\vX^{\underset{\bullet}{1}}$ and $\vX^{\underset{\bullet}{2}}$ for 
nodes ${\underset{\bullet}{1}}$ and ${\underset{\bullet}{2}}$, respectively.
These nodes are moved into the current configuration 
at $\vx^{\underset{\bullet}{1}}$ and $\vx^{\underset{\bullet}{2}}$ through 
nodal displacements $\vu^{\underset{\bullet}{1}}$ and
$\vu^{\underset{\bullet}{2}}$.  

%\begin{figure*}
%\begin{figure}[htb]
% \psfrag{A}{$\vE_1$}
% \psfrag{B}{$\vE_2$}
% \psfrag{C}{$\vE_3$}    
% \psfrag{O}{$O$}
% \psfrag{L}{$L$}
% \psfrag{l}{$\ell$}
% \psfrag{X}{$\vX^{\underset{\bullet}{1}}$}  
% \psfrag{x}{$\vx^{\underset{\bullet}{1}}$}  
% \psfrag{Y}{$\vX^{\underset{\bullet}{2}}$}  
% \psfrag{y}{$\vx^{\underset{\bullet}{2}}$}        
% \psfrag{T}{$\vThat$}
% \psfrag{t}{$\vthat$}
% %\fbox{\begin{minipage}{0.98\textwidth}\centering
% \begin{center}
% \includegraphics[width=0.8\textwidth]{fig/nonlinearTruss.eps}
% \end{center}
% \caption{Nonlinear truss in reference and current configurations.}
% \label{fig:nonlinearTruss} % label must come after caption 
% %\end{minipage}
% %} % \fbox
%%\end{figure*} 
%\end{figure}

\begin{figure}[htb]
  \begin{center}
  %\begin{overpic}[width=\textwidth, grid, tics=5]{fig/nonlinear_truss}
  %\begin{overpic}[width=\textwidth,natwidth=371,natheight=238]{fig/nonlinear_truss.pdf}
  \begin{overpic}[width=\textwidth]{fig/nonlinear_truss.pdf}
 \put(19,  8){$\vE_1$}
 \put( 7, 19){$\vE_2$}
 \put( 0,  2){$\vE_3$}    
 \put( 9, 10){$O$}
 \put(15.7, 48.9){$L$}
 \put(67.6, 32){$\ell$}
 \put( 4, 42){$\vX^{\underset{\bullet}{1}}$}  
 \put(47, 27){$\vx^{\underset{\bullet}{1}}$}  
 \put(25, 60){$\vX^{\underset{\bullet}{2}}$}  
 \put(82, 40){$\vx^{\underset{\bullet}{2}}$}        
 \put(39, 60){$\vThat$}
 \put(97, 37){$\vthat$}
  \end{overpic}
  \end{center}
 \caption{Nonlinear truss in reference and current configurations.}
 \label{fig:nonlinear_truss} % label must come after caption
\end{figure}

The strain energy of the truss is the strain energy 
density integrated over the truss's volume, 
\be
 \cW  = \int_{\Omega} \vsigma : \vhstrain \; d\Omega.
\ee
The virtual work done by the truss is
\be
 \delta \cW = \int_{\Omega} \delta \vhstrain : \vsigma \; d\Omega.
\ee
The virtual work computed in the reference configuration is then 
\begin{align}
  \delta \cW & = \int_{L} \delta \vhstrain : \vsigma A \; dL, \\
  & \overset{\mbox{\footnotesize (1D)}}{=} 
  \delta \left(\dst\frac{\ell - L}{L} \right) \sigma A L 
  = \dst\frac{\delta \ell}{L} \sigma A L = \delta \ell \sigma A, \\
  & \overset{(2.70)}{=} 
  (- \delta \vu^{\underset{\bullet}{1}} + 
  \delta \vu^{\underset{\bullet}{2}} ) \cdot \vthat 
  \; \sigma A, \label{eq:NLTrussFirstVariation} \\
  & \overset{\mbox{\footnotesize (lin)}}{=} 
  (- \delta \vu^{\underset{\bullet}{1}} + 
  \delta \vu^{\underset{\bullet}{2}} ) \cdot \vthat  \; 
  E \left( \dst\frac{\ell - L}{L} \right) A, \\
  & = (- \delta \vu^{\underset{\bullet}{1}} + 
  \delta \vu^{\underset{\bullet}{2}} ) \cdot \vthat \; k (\ell - L),
\end{align}
where the result from 
Hovey\footnote{Hovey CB (2000). Stress and Motion in Diarthrodial 
Joints Using Coupled Nonlinear Elastodynamics, Rigid Body Dynamics, 
and Large-Slip Contact-Impact, Doctoral Thesis, Stanford University, 
Stanford, CA.}~page~45, Eq.~(2.70), 
$\delta \ell = (- \delta \vu^{\underset{\bullet}{1}} + 
\delta \vu^{\underset{\bullet}{2}} ) \cdot \vthat$, has been used, 
and agrees with Eq.~(2.76) of the same manuscript.
The virtual work expressed in terms of nodal displacements and the 
{\em nonlinear} internal force vector in~(\ref{eq:NLTrussFirstVariation}) is then
\begin{equation}
 \boxed{
  \delta \cW = \vdelta \cdot \vfint = 
   \left\{\begin{tabular}{c}
    $\delta \vu^{\underset{\bullet}{1}}$ \\ $\delta \vu^{\underset{\bullet}{2}}$
   \end{tabular}\right\}^T
   \sigma A 
   \left\{
	  \begin{tabular}{c}
	    $-\vthat$ \\ $\vthat$
          \end{tabular}
   \right\}
   } \; .  
   \label{eq:NLTrussFirstVariationFint}
\end{equation}
With $\delta \ell$ at hand, and noting that
\begin{align}
 \delta \vthat & = \delta \left(
 \frac{\vx_2 - \vx_1}{l}\right) = \delta 
 \left((\vx_2 - \vx_1)l^{-1} \right), \\
  % & = (\delta \vu_2 - \delta \vu_1) l^{-1} + (\vx_2 - \vx_1)(-l^{-2}) 
  %  \delta l, \\
  % & = \frac{\delta \vu_2 - \delta \vu_1}{l} - 
  %   \frac{\vx_2 - \vx_1}{l^2} (-\delta \vu_1 + \delta \vu_2) \cdot \vthat, \\
  & = \frac{1}{l}\left[\vid - \vthat \otimes \vthat
  \right] (-\delta \vu_1 + \delta \vu_2),
\end{align} 
the second variation of the nonlinear form in~(\ref{eq:NLTrussFirstVariation}), 
or equivalently in~(\ref{eq:NLTrussFirstVariationFint}), 
is found through
%\clearpage
\begin{align}
  \Delta (\delta \cW) 
   & = \Delta \left[(- \delta \vu^{\underset{\bullet}{1}} + 
   \delta \vu^{\underset{\bullet}{2}} ) \cdot \sigma A \; \vthat \right], \\
   % & = (- \delta \vu^{\underset{\bullet}{1}} + \delta \vu^{\underset{\bullet}{2}} ) 
   % \cdot \Delta \vthat \; \sigma A \;\;+\;\; (- \delta \vu^{\underset{\bullet}{1}} + \delta 
   % \vu^{\underset{\bullet}{2}} ) \cdot \vthat \; \Delta \sigma A, \\
   & = (- \delta \vu^{\underset{\bullet}{1}} 
   + \delta \vu^{\underset{\bullet}{2}} ) 
   \cdot \Delta \sigma A \; \vthat \nonumber \\
   & \hspace{0.5cm} + (- \delta \vu^{\underset{\bullet}{1}} + \delta \vu^{\underset{\bullet}{2}} ) \cdot \sigma A \; \Delta \vthat.
\end{align}
Now, adopting a compact notation where 
$\delta \vu = 
(- \delta \vu^{\underset{\bullet}{1}} + \delta \vu^{\underset{\bullet}{2}} )$ 
and 
$\Delta \vu = (- \Delta \vu^{\underset{\bullet}{1}} + \Delta \vu^{\underset{\bullet}{2}} )$,
\begin{align}
 &\Delta (\delta \cW) \nonumber \\
 &= \delta \vu \cdot 
 \left\{ 
  \dst\frac{\partial \sigma}{\partial \hstrain} \Delta \hstrain A \; \vthat     
   +
  \sigma A \frac{1}{\ell}\left[\vid - \vthat \otimes \vthat
 \right] \Delta \vu
 \right\}, \\
 % & = (- \delta \vu^{\underset{\bullet}{1}} + \delta \vu^{\underset{\bullet}{2}} ) 
 % \cdot \left\{ \frac{1}{\ell}\left[\vid - \vthat \otimes \vthat
 %   \right] (- \Delta \vu^{\underset{\bullet}{1}} + \Delta \vu^{\underset{\bullet}{2}}) \; 
 %\sigma A \;\;+\;\; \vthat \; \dst\frac{\partial \sigma}{\partial \hstrain} \Delta \hstrain A
 %\right\}, \\
 & = \delta \vu \cdot 
 \left\{ 
   \dst\frac{\partial \sigma}{\partial \hstrain} 
     \left(\dst\frac{\Delta \ell}{L} \right)
     A \; \vthat 
    + 
    \frac{\sigma A}{\ell}\left[\vid - \vthat \otimes \vthat
    \right] \Delta \vu \right\}.
\end{align}
Expanding $\delta \vu$ and $\Delta \vu$ gives
\begin{align}
\Delta(\delta \cW) & = (- \delta \vu^{\underset{\bullet}{1}} + \delta \vu^{\underset{\bullet}{2}} ) \nonumber \\
 &
 \hspace{0.5cm} 
  \cdot 
   \left[ 
    \dst\frac{\partial \sigma}{\partial \hstrain} \dst\frac{A}{L} \;
    \vthat \otimes \vthat 
    \;\;+\;\; 
    \frac{\sigma A}{\ell}\left[\vid - \vthat \otimes \vthat\right]  
   \right] \nonumber \\
 %
  &  \hspace{1.0cm} \left\{ 
  (- \Delta \vu^{\underset{\bullet}{1}} + \Delta \vu^{\underset{\bullet}{2}})
  \right\}.
\end{align}
This leads to the element stiffness equations
\begin{align}
  \Delta \delta \cW &= \nonumber \\
  & \hspace{2.5\columnwidth} \nonumber
\end{align}
  \vspace{-1cm}
\begin{equation}
  \boxed{
  \vdelta \cdot \vK \vDelta
   = \left\{
	\begin{tabular}{c}
        $\delta \vu^{\underset{\bullet}{1}}$ \\ 
	$\delta \vu^{\underset{\bullet}{2}}$
	\end{tabular}
	\right\}^T
	\left[
	 \begin{tabular}{rr}
	   	$\vkappa$  & $-\vkappa$ \\ 
	 	$-\vkappa$ & $\vkappa$
	 \end{tabular}
	\right]
	\left\{
	 \begin{tabular}{c}
	  	$\Delta \vu^{\underset{\bullet}{1}}$ \\ 
		$\Delta \vu^{\underset{\bullet}{2}}$
	  \end{tabular}
	 \right\}
	 } \; ,
\end{equation}
where
\begin{equation}
 \vkappa = \left[\dst\frac{\partial \sigma}{\partial \hstrain}
 \dst\frac{A}{L} \; \vthat \otimes \vthat + \frac{\sigma A}{\ell}
 \left[\vid - \vthat \otimes \vthat \right] \right].
\end{equation}
Note that the explicit form of $\partial \sigma / \partial \hstrain$ 
will depend on the plastic multiplier $\gammadot$.  If $\gammadot = 0$, 
then $\partial \sigma / \partial \hstrain = E$, as shown in the 
elastic region of Fig.~\ref{fig:stress_strain_harden}~(a).  
If $\gammadot > 0$, then $\partial \sigma / \partial \hstrain = \Eep$, as 
shown in the elasto-plastic region of Fig.~\ref{fig:stress_strain_harden}~(a).  


% ========================
\section{From 1D to 3D}
In one-dimensional plasticity, we spoke of the Cauchy stress scalar
$\sigma \in \real$ with positive values describing tension (resulting
in elongations) and compression (resulting in contractions).  
Now, we generalize this concept to speak of the Cauchy stress
tensor $\vsigma$ as a collection of all the normal and shear 
Cauchy stress components present in three dimensions.  

We decompose $\vsigma$ into volumetric and deviatoric components, 
denoted
\be
\vsigma = \vsigmavol + \vsigmadev.
\ee
Next, we {\em define} each of these two components.  
The {\bf volumetric Cauchy stress tensor} is defined as
\be
\vsigmavol \defe \frac{1}{3} \trace(\vsigma) \vid = 
\frac{1}{3} \sigma_{kk} \; \delta_{ij}.
\ee
It is convenient to define the mean stress, $\bar{\sigma}$, as
\be
\bar{\sigma} \defe \frac{1}{3} 
(\sigma_{11} + \sigma_{22} + \sigma_{33}).  
\ee
Then 
\be
\vsigmavol \defe \bar{\sigma} \vid = 
\bar{\sigma} \; \delta_{ij}.
\ee
The {\bf deviatoric Cauchy stress tensor},  
$\tsbold$, is defined as
\be
\tsbold \defe \vsigma - \vsigmavol = \vsigma - \bar{\sigma} \; \vid.
\ee 
It is often convenient to define {\bf pressure} $P \in \real$ as
\be
P \defe - \frac{1}{3} \trace(\vsigma) \hspace{1.0cm}
  \left\{
    \begin{tabular}{l}
      $P > 0$ compression, \\
      $P < 0$ tension.
    \end{tabular}
  \right.
\ee
Note the negative sign in the definition.  The negative sign
allows us to harmonize our concepts of (compressive, e.g., 
underwater) pressure as a positive, scalar with 
tension/compression (positive/negative) convention of the Cauchy
stress tensor.


\subsection{Uniaxial Tension Test}
With these definitions at hand, we are now poised to demonstrate
how the results of a uniaxial tension test at the plastic limit
are used to generalize a theory of small-strain plasticity to
three dimensions. 

A yielding, the Cauchy stress tensor $\vsigma$ has the simple form
\be
\left.
\vsigma
\right|_y = 
\left[
  \begin{tabular}{lll}
    $\sigma_y$   & 0 & 0 \\
                 & 0 & 0 \\
     &   & 0 
  \end{tabular}
\right].
\ee

\begin{Hremark}
Note that for symmetric matrices, we will 
avoid filling in the lower part of the matrix both for brevity, and to
reinforce the idea of symmetry.
\end{Hremark}

The volumetric part of the Cauchy stress tensor at yield is then
\be
\left.
\vsigmavol 
\right|_y
= \frac{\sigma_y}{3}
\left[
  \begin{tabular}{lll}
    1            & 0 & 0 \\
                 & 1 & 0 \\
     &   & 1 
  \end{tabular}
\right].
\ee
The deviatoric part of the Cauchy stress tensor at yield is then
\be
\left.
\tsbold 
\right|_y
= \frac{3 \;\sigma_y}{3}
\left[
  \begin{tabular}{lll}
    1            & 0 & 0 \\
                 & 1 & 0 \\
     &   & 1 
  \end{tabular}
\right]
 - 
\frac{\sigma_y}{3}
\left[
  \begin{tabular}{lll}
    1            & 0 & 0 \\
                 & 1 & 0 \\
     &   & 1 
  \end{tabular}
\right] 
 = 
\frac{\sigma_y}{3}
\left[
  \begin{tabular}{lll}
    2 & \;0 & \;0 \\
      &-1 & \;0 \\
      &   &-1 
  \end{tabular}
\right].
\label{eq:deviatoric_uniaxial_tension}
\ee
So, we now have the form of the deviatoric Cauchy stress
tensor $\tsbold$ at yield.  How can this result for a uniaxial tension test
be generalized for three dimensions?  The thought experiment is to
envision $\tsbold$ for {\em any and all} matrix forms, not just the form
obtained in (\ref{eq:deviatoric_uniaxial_tension}).  Any and all matrix
form for $\tsbold$ will have one commonality:  their matrix invariants.

\subsubsection{The Cauchy Stress Tensor Invariants}
The invariants of the Cauchy stress tensor, $\vsigma$, satisfy the
characteristic equation (see (\ref{eq:invariants_sot}) for details),
\be
0 = - \lambda^3 + I_1 \lambda^2 - I_2 \lambda^1 + I_3 \lambda^0,
\ee
where the three invariants, $I_1$, $I_2$, and $I_3$ are 
\begin{align}
 I_1 &= \trace(\vsigma), \\
 I_2 &= \half\left[ \trace(\vsigma) \trace(\vsigma) 
        - \vsigma : \vsigma \right] 
	 = \half\left[I^2_1 - \vsigma : \vsigma \right], \\
 I_3 &= \det(\vsigma).
\end{align}

\subsubsection{The Deviatoric Cauchy Stress Tensor Invariants}
The invariants of the deviatoric Cauchy stress tensor, $\tsbold$, 
satisfy the characteristic equation
\be
0 = - \lambda^3 + J_1 \lambda^2 - J_2 \lambda^1 + J_3 \lambda^0,
\ee
where the three invariants, $J_1$, $J_2$, and $J_3$ are
\begin{align}
 J_1 &= \trace(\tsbold) = 0, \\
 J_2 &= \half\left[ \trace(\tsbold) \trace(\tsbold) 
        - \tsbold : \tsbold \right] 
	 = \half\left[J^2_1 - \tsbold : \tsbold \right] 
	 = -\half \tsbold : \tsbold, \\
 J_3 &= \det(\tsbold).
\end{align}


will have a first invariant, 
$J_1$ equal to zero since the trace of a deviatoric second order tensor
is zero.  


Then the first invariant of $\tsbold$, 
\be
J_1(\tsbold) \defe \trace(\tsbold) = s_{kk} = 0,
\ee
is uninteresting because it has
zero value for all stress states of $\vsigma$.
The second invariant of $\tsbold$, 
\be
J_2(\tsbold) \defe \half \left( 
              \trace \left( \tsbold \cdot \tsbold \right) \right) 
              = \half s_{ij} s_{ji} 
\ee
is interesting because we arrive at a quadratic form of $\tsbold$.  
Returning to the thought experiment, we can now envision of 
form of $\tsbold$ by way of $J_2(\tsbold)$ as valid for all 
values Cauchy stress tensors $\vsigma$ and with plasticity occurring
when 
$J_2(\tsbold)$ reaches $J_2(\left. \tsbold \right|_y)$.
Said equivalently, we can imagine a yield function, 
as defined for scalars in (\ref{eq:yieldFunctionPerfectPlastic}), recast
now for the entire stress tensor $\vsigma$ as
\be
 \Phi(\vsigma) \defe J_2(\tsbold(\vsigma)) - J_2(\left. \tsbold \right|_y)
 \leq 0.
\ee
Now 
\be
J_2(\left. \tsbold \right|_y) 
= \half \; \frac{\sigma_y^2}{9} \; (4 + 1 + 1) = \frac{\sigma_y^2}{3}.
\ee
Thus
\be
 \Phi(\vsigma) \defe J_2(\tsbold(\vsigma)) - 
 \frac{\sigma_y^2}{3} 
 \leq 0.
\ee
Now, on the yield surface, $\Phi = 0$, which means
\be
 J_2(\tsbold(\left. \vsigma \right|_y)) 
 = 
 \frac{\sigma_y^2}{3} 
 \implies 
 \sqrt{ 3 J_2(\tsbold(\left. \vsigma \right|_y)) }
 = 
 \sigma_y.
\ee
This gives rise to the form of the yield function of
so-called {\bf $J_2$ plasticity},
\be
\boxed{
 \Phi(\vsigma) \defe \sqrt{ 3 J_2 (\tsbold(\vsigma)) } - \sigma_y 
 \leq 0.
}
\ee
The quantity $\sqrt{ 3 J_2 }$ is often referred to
as the {\bf von Mises stress} or {\bf equivalent tensile stress}
and denoted as $\sigma_{\mbox{\tiny VM}}$.
Thus, the yield function may be equivalently written as
\be
\boxed{
 \Phi(\vsigma) \defe \sigma_{\mbox{\tiny VM}}(\vsigma) - \sigma_y 
 \leq 0.
}
\ee

\begin{Hresult}
The squared von Mises stress can be written explicity in terms of the Cauchy stress components as follows:
\be
\sigma^2_{\mbox{\tiny VM}} = 3 J_2 = \frac{(\sigma_{11} - \sigma_{22})^2 + (\sigma_{22} - \sigma_{33})^2 + (\sigma_{33} - \sigma_{11})^2}{2} + 3 (\sigma^2_{12} + \sigma^2_{23} + \sigma^2_{31}).
\ee
\end{Hresult}

\begin{Hproof}
Recall $J_2(\ts) = \half\;\tsbold : \tsbold = \half \;\trace( \tsbold \cdot \tsbold)$.  Also, $\tsbold \defe \vsigma - \bar{\sigma} \vid$, where $\bar{\sigma} \defe \frac{1}{3}(\sigma_{11} + \sigma_{22} + \sigma_{33})$.
Then, 
\begin{align}
&\sigma^2_{\mbox{\tiny VM}} = 3 J_2 \\
&= \frac{3}{2} 
\trace 
\left(
\left[
\begin{tabular}{ccc}
  $(\sigma_{11} - \bar{\sigma})$ & $\sigma_{12}$ & $\sigma_{13}$ \\
  & $(\sigma_{22} - \bar{\sigma})$ & $\sigma_{23}$ \\
  sym. & & $(\sigma_{33} - \bar{\sigma})$ 
\end{tabular}
\right]
\left[
\begin{tabular}{ccc}
  $(\sigma_{11} - \bar{\sigma})$ & $\sigma_{12}$ & $\sigma_{13}$ \\
  & $(\sigma_{22} - \bar{\sigma})$ & $\sigma_{23}$ \\
  sym. & & $(\sigma_{33} - \bar{\sigma})$ 
\end{tabular}
\right]
\right),
\\
% \end{align}
% \begin{align}
&= \frac{3}{2}\left[
   (\sigma_{11} - \bar{\sigma})^2 + 
   (\sigma_{22} - \bar{\sigma})^2 + 
   (\sigma_{33} - \bar{\sigma})^2 +
   2 \sigma^2_{12} + 
   2 \sigma^2_{23} + 
   2 \sigma^2_{31}
\right],
\\
&=
\frac{3}{2}\left[
  \sigma^2_{11} + \sigma^2_{22} + \sigma^2_{33} - 
  2 \bar{\sigma} (\sigma_{11} + \sigma_{22} + \sigma_{33}) + 
  3 \bar{\sigma}^2
\right]
+ 3 (\sigma^2_{12} + \sigma^2_{23} + \sigma^2_{31}),
\\
&= 
\frac{3}{2}
\left[
  \sigma^2_{11} + \sigma^2_{22} + \sigma^2_{33} - 
  \frac{2}{3}(\sigma_{11} + \sigma_{22} + \sigma_{33})^2 +
  \frac{3}{9}(\sigma_{11} + \sigma_{22} + \sigma_{33})^2
\right]
+ 3 (\sigma^2_{12} + \sigma^2_{23} + \sigma^2_{31}),
\\
&=
\half
\left[
3 (\sigma^2_{11} + \sigma^2_{22} + \sigma^2_{33}) 
- (\sigma_{11} + \sigma_{22} + \sigma_{33})^2
\right]
+ 3 (\sigma^2_{12} + \sigma^2_{23} + \sigma^2_{31}),
\\
&=
\half
\left[
2 (\sigma^2_{11} + \sigma^2_{22} + \sigma^2_{33} 
- \sigma_{11} \sigma_{22} 
- \sigma_{22} \sigma_{33}
- \sigma_{33} \sigma_{11}
)
\right]
+ 3 (\sigma^2_{12} + \sigma^2_{23} + \sigma^2_{31}),
\\
&=
\frac{
(\sigma_{11} - \sigma_{22})^2 + 
(\sigma_{22} - \sigma_{33})^2 + 
(\sigma_{33} - \sigma_{11})^2}{2} 
+ 3 (\sigma^2_{12} + \sigma^2_{23} + \sigma^2_{31}),
\end{align}
which is the result.
\end{Hproof}

